\section{Scaling Laws and Their Limits}

\begin{frame}{}
    \LARGE Next-Generation Architectures: \textbf{Scaling Laws and Their Limits}

    \vspace{1em}
    \normalsize Given a fixed budget, how big should the model be, and how long should it train?
\end{frame}

\begin{frame}{The Textbook Story, and What 2024 Corrected}
    \begin{itemize}\itemsep0.5em
        \item \textbf{Recap.} Loss follows a power law in parameters $N$ and tokens $D$. For a compute budget $C\approx 6ND$ there is an optimal split.
        \begin{itemize}\itemsep0.2em
            \item Kaplan et al. (2020): grow the model faster, $N^*\propto C^{0.73}$.
            \item Chinchilla (2022): grow both equally, $N^*\propto C^{0.5}$. About \textbf{20 tokens per parameter}.
        \end{itemize}
        \item \textbf{Correction 1.} Chinchilla's own parametric fit was wrong. Taken literally it implies about 70 tokens per parameter. A careful refit recovers 20, and shows the reported confidence intervals were implausibly tight.
        \item \textbf{Correction 2.} Kaplan and Chinchilla never disagreed about Transformers. The gap is a \emph{measurement artefact}, as the next slide shows.
    \end{itemize}
    \vspace{0.2em}
    \takeaway{``20 tokens per parameter'' is a property of one training setup. It is a useful heuristic, not a law.}
    \src{Hoffmann et al., \href{https://arxiv.org/abs/2203.15556}{arXiv:2203.15556}. Besiroglu et al., ``Chinchilla Scaling: A replication attempt,'' \href{https://arxiv.org/abs/2404.10102}{arXiv:2404.10102}. Pearce and Song, \href{https://arxiv.org/abs/2406.12907}{arXiv:2406.12907}.}
\end{frame}

\begin{frame}{From Kaplan to Chinchilla in Four Fixes}
    \begin{columns}[T,onlytextwidth]
        \column{0.5\textwidth}
        \centering
        \begin{adjustbox}{max width=\linewidth}\begin{tikzpicture}
            \begin{axis}[ybar, width=7.2cm, height=5.6cm, bar width=0.55cm,
                ymin=0.4, ymax=0.95, ylabel={\scriptsize exponent $a$ in $N^*\propto C^{a}$},
                symbolic x coords={A,B,C,D,E}, xtick=data,
                xticklabels={Kaplan setup, + head FLOPs, + scaled warmup, + cosine decay, + tuned optimiser},
                x tick label style={font=\tiny, rotate=25, anchor=east},
                y tick label style={font=\scriptsize},
                nodes near coords, nodes near coords style={font=\scriptsize},
                axis lines=left, enlarge x limits=0.15]
                \addplot[fill=KAUSTcyan!55, draw=KAUSTcyan!70!black] coordinates {(A,0.88) (B,0.70) (C,0.60) (D,0.57) (E,0.50)};
            \end{axis}
        \end{tikzpicture}\end{adjustbox}
        \column{0.48\textwidth}
        \begin{itemize}\itemsep0.5em
            \item Over 900 training runs reproduce Kaplan's exponent, then remove it step by step.
            \item The biggest single cause: \textbf{not counting the compute of the output layer}, which matters at small scale.
            \item A side result against Chinchilla: learning-rate decay is \emph{not} needed for the 0.5 law.
            \item \textbf{Lesson for your own fits:} scaling exponents are fragile to protocol. Small models and untuned hyperparameters bias them.
        \end{itemize}
    \end{columns}
    \src{Porian et al., ``Resolving Discrepancies in Compute-Optimal Scaling of Language Models,'' \href{https://arxiv.org/abs/2406.19146}{arXiv:2406.19146}, Figure 1.}
\end{frame}

\begin{frame}{Compute-Optimal Is Not Deployment-Optimal}
    \begin{columns}[T,onlytextwidth]
        \column{0.4\textwidth}
        \centering
        \small
        \begin{adjustbox}{max width=\linewidth}\begin{tabular}{lr}
            \toprule
            \textbf{Model} & \textbf{Tokens / param} \\
            \midrule
            GPT-3 & 2 \\
            Chinchilla & 20 \\
            Llama 2 70B & 29 \\
            Mistral 7B & 110 \\
            Llama 3 70B & 215 \\
            \rowcolor{KAUSTcyan!12} Llama 3 8B & about 1,900 \\
            \bottomrule
        \end{tabular}\end{adjustbox}

        \vspace{0.3em}
        {\scriptsize Last row derived from 15T tokens and 8B parameters.}
        \column{0.58\textwidth}
        \begin{itemize}\itemsep0.5em
            \item Chinchilla minimises \textbf{training} compute for a target loss. A deployed model also pays for every token it ever generates.
            \item If inference demand is large, the best choice is a \textbf{smaller model trained far longer}. Quality keeps improving up to 10,000 tokens per parameter.
            \item Meta's own account: the Chinchilla-optimal budget for an 8B model is about 200B tokens. Llama 3 8B kept improving log-linearly up to 15T.
            \item Flagships stay near compute-optimal. \textbf{Small models are over-trained by 100$\times$ or more.}
        \end{itemize}
    \end{columns}
    \vspace{0.3em}
    \src{Table: Hashimoto, Stanford CS336 Lecture 9 (2026). Sardana et al., ``Beyond Chinchilla-Optimal,'' \href{https://arxiv.org/abs/2401.00448}{arXiv:2401.00448}. Meta, ``Introducing Meta Llama 3'' (Apr 2024).}
\end{frame}

\begin{frame}{When the Data Runs Out}
    \begin{itemize}\itemsep0.33em
        \item Scaling laws assume fresh data. What if you must repeat it?
        \item \textbf{Repetition has a budget.} Up to about 4 epochs is almost as good as new data. The value of further epochs decays exponentially:
        \[
            D' = U_D + U_D\,R_D^*\big(1-e^{-R_D/R_D^*}\big), \qquad R_D^*\approx 15
        \]
        where $U_D$ is unique tokens and $R_D$ the number of repetitions.
        \item \textbf{The stock is finite.} Epoch AI estimates about 300T tokens of usable public text, fully used some time between 2026 and 2032.
        \item \textbf{Responses:} stronger regularisation, synthetic data for maths and code, and compute spent on reinforcement learning instead of more text.
    \end{itemize}
    \vspace{0.07em}
    \takeaway{Sutskever, NeurIPS 2024: ``We've achieved peak data and there'll be no more. There's only one internet.''}
    \src{Muennighoff et al., ``Scaling Data-Constrained Language Models,'' \href{https://arxiv.org/abs/2305.16264}{arXiv:2305.16264}. Villalobos et al., ``Will we run out of data?'' Epoch AI (2024). Kim et al., \href{https://arxiv.org/abs/2509.14786}{arXiv:2509.14786}. Quote as reported by The Verge (13 Dec 2024).}
\end{frame}

\begin{frame}{Data Quality Is a Compute Multiplier}
    \begin{center}
    \small
    \renewcommand{\arraystretch}{1.05}
    \begin{adjustbox}{max width=\linewidth}\begin{tabular}{>{\raggedright\arraybackslash}p{2.6cm}>{\raggedright\arraybackslash}p{5.2cm}>{\raggedright\arraybackslash}p{5.6cm}}
        \toprule
        \textbf{Study} & \textbf{What was done} & \textbf{Effect} \\
        \midrule
        FineWeb-Edu (2024) & keep web pages an LLM rates as educational & MMLU 33 to 37\%, ARC 46 to 57\%. Matches the baseline with about 10$\times$ fewer tokens \\
        DCLM (2024) & a simple fastText quality classifier over a 240T-token pool & comparable to Llama 3 8B on MMLU with about 6.6$\times$ less compute \\
        Data mixing laws (2025) & predict loss as a function of mixture proportions & the optimised mix matches the default mix trained 48\% longer \\
        \bottomrule
    \end{tabular}\end{adjustbox}
    \end{center}
    \begin{itemize}\itemsep0.29em
        \item Classical scaling laws have no term for quality. In practice it shifts the whole curve.
        \item \textbf{The tension:} aggressive filtering shrinks the pool of unique tokens, which brings back the repetition problem of the previous slide.
        \item The best filter depends on the budget. With more compute, filter less.
    \end{itemize}
    \src{Penedo et al., ``The FineWeb Datasets,'' \href{https://arxiv.org/abs/2406.17557}{arXiv:2406.17557}. Li et al., ``DataComp-LM,'' \href{https://arxiv.org/abs/2406.11794}{arXiv:2406.11794}. Ye et al., ``Data Mixing Laws,'' \href{https://arxiv.org/abs/2403.16952}{arXiv:2403.16952}. Goyal et al., \href{https://arxiv.org/abs/2404.07177}{arXiv:2404.07177}.}
\end{frame}

\begin{frame}{Does Model Shape Matter?}
    \begin{itemize}\itemsep0.4em
        \item Kaplan's finding still holds at scale: for a fixed parameter count, loss depends only weakly on shape. So shape is set by convention and by hardware.
    \end{itemize}
    \begin{center}
    \small
    \begin{adjustbox}{max width=\linewidth}\begin{tabular}{lll}
        \toprule
        \textbf{Choice} & \textbf{Convention} & \textbf{Recent outliers} \\
        \midrule
        FFN width / model width & 4, or 8/3 with gated units & Gemma 2 uses 8 \\
        Model width / depth & 100 to 200 & Gemma 3: 87. Gemma 4: 61 \\
        Heads $\times$ head dim / width & 1 & Qwen 3.5: 1.2 \\
        Vocabulary & 100K to 250K in production & Gemma 4: 262K \\
        \bottomrule
    \end{tabular}\end{adjustbox}
    \end{center}
    \begin{itemize}\itemsep0.4em
        \item \textbf{Where shape does matter:}
        \begin{itemize}\itemsep0.2em
            \item \emph{Small models.} Below a billion parameters, deep and thin beats shallow and wide (MobileLLM).
            \item \emph{Vocabulary.} The optimal size grows with compute. By one estimate, Llama 2 70B's 32K vocabulary should have been at least 216K.
        \end{itemize}
    \end{itemize}
    \src{Table: Hashimoto, Stanford CS336 Lecture 3 (2026). Liu et al., ``MobileLLM,'' \href{https://arxiv.org/abs/2402.14905}{arXiv:2402.14905}. Tao et al., ``Scaling Laws with Vocabulary,'' \href{https://arxiv.org/abs/2407.13623}{arXiv:2407.13623}.}
\end{frame}

\begin{frame}{Scaling Laws Grew New Axes}
    \begin{itemize}\itemsep0.39em
        \item \textbf{Precision.} Treat bit-width as a third resource. Low precision acts like a reduction in effective parameters:
        \[
            N_{\text{eff}} = N\,\big(1-e^{-P_w/\gamma_w}\big)\big(1-e^{-P_a/\gamma_a}\big)\big(1-e^{-P_{kv}/\gamma_{kv}}\big)
        \]
        Two consequences. About 7 to 8 bits is compute-optimal for training. And \alert{over-trained models quantise worse}, which is exactly the small deployed models of two slides ago.
        \item \textbf{Architecture.} Hybrids, MoE and multimodal models get their own fitted laws. The exponents stay close to the dense Transformer's.
        \item \textbf{Emergence.} Are sudden jumps in ability real? Two answers that fit together:
        \begin{itemize}\itemsep0.13em
            \item with continuous metrics, most ``emergent'' curves become smooth
            \item abilities appear when the \emph{pre-training loss} crosses a threshold, whatever the model size
        \end{itemize}
        \item So the working recipe is two-stage: predict loss from compute, then task score from loss.
    \end{itemize}
    \src{Kumar et al., ``Scaling Laws for Precision,'' \href{https://arxiv.org/abs/2411.04330}{arXiv:2411.04330}. Schaeffer et al., \href{https://arxiv.org/abs/2304.15004}{arXiv:2304.15004}. Du et al., \href{https://arxiv.org/abs/2403.15796}{arXiv:2403.15796}. MiniMax, \href{https://arxiv.org/abs/2501.08313}{arXiv:2501.08313}.}
\end{frame}

\begin{frame}{A Third Place to Spend Compute: Inference}
    \begin{itemize}\itemsep0.55em
        \item Training laws trade training FLOPs for loss. Since 2024 there is a second curve: \textbf{inference FLOPs for task success}.
        \item OpenAI's o1 on AIME 2024 mathematics problems:
    \end{itemize}
    \begin{center}
    \small
    \begin{adjustbox}{max width=\linewidth}\begin{tabular}{lc}
        \toprule
        \textbf{Setting} & \textbf{Problems solved} \\
        \midrule
        GPT-4o & 12\% \\
        o1, one sample & 74\% \\
        o1, majority vote over 64 samples & 83\% \\
        o1, 1,000 samples re-ranked by a learned scorer & 93\% \\
        \bottomrule
    \end{tabular}\end{adjustbox}
    \end{center}
    \begin{itemize}\itemsep0.55em
        \item At matched FLOPs, a small model with well-allocated test-time compute can beat a model \textbf{14$\times$ larger}.
        \item The catch: gains depend on a \textbf{verifier}. Without one, voting saturates after a few hundred samples.
    \end{itemize}
    \src{OpenAI, ``Learning to reason with LLMs'' (Sep 2024). Snell et al., \href{https://arxiv.org/abs/2408.03314}{arXiv:2408.03314}. Brown et al., ``Large Language Monkeys,'' \href{https://arxiv.org/abs/2407.21787}{arXiv:2407.21787}.}
\end{frame}

\begin{frame}{Limits: Where Scaling Stands in October 2026}
    \begin{itemize}\itemsep0.36em
        \item \textbf{Compute keeps growing.} Frontier training compute has risen about 5$\times$ per year since 2020. The largest estimated run so far is Grok 4, at about $5\times10^{26}$ FLOP.
        \item \textbf{What binds next is power, then data.} Epoch AI projects runs near $2\times10^{29}$ FLOP as feasible by 2030, with electricity the first hard ceiling.
        \item \textbf{Pre-training alone pays less.} GPT-4.5 used roughly ten times GPT-4's compute for incremental gains. GPT-5 then used \emph{less} training compute than GPT-4.5: scaling post-training on a smaller model gave better returns.
        \item \textbf{Reinforcement learning has its own curve.} It is sigmoidal in compute, not a power law. Most recipe choices change how fast you approach the ceiling. Few change the ceiling.
    \end{itemize}
    \vspace{0.13em}
    \caveat{Be careful with ``scaling has hit a wall''. No lab has published a frontier-scale loss curve that shows the power law breaking.}
    \src{Epoch AI, trends dashboard (Feb 2026), ``Can AI scaling continue through 2030?'' (2024) and ``Why GPT-5 used less training compute than GPT-4.5'' (Sep 2025). Lambert, Interconnects (Feb 2025). Meta, ``The Art of Scaling RL Compute for LLMs,'' \href{https://arxiv.org/abs/2510.13786}{arXiv:2510.13786}.}
\end{frame}

\begin{frame}{Scaling: What to Remember}
    \begin{itemize}\itemsep0.6em
        \item The compute-optimal rule survives, but \textbf{its constants depend on protocol}. Fit your own.
        \item Deployed models are deliberately \textbf{over-trained}, because inference dominates lifetime cost.
        \item \textbf{Data} is now a binding constraint: its quantity, its quality and how often it can be repeated.
        \item Compute has three destinations: \textbf{pre-training, post-training and inference}. Each has a different curve.
    \end{itemize}
    \vspace{0.3em}
    \vspace{0.3em}
    \takeaway{\textbf{Next.} All of this assumed text. The last section asks what changes when the model must also see, hear and draw.}
\end{frame}
